% GENERATED FILE. Edit canonical lessons, then run npm run generate:books.
% canonical-source-hash: fnv1a64:2255f0215a750963
% canonical-lessons: ES-C764-inolvidable, ES-C764-olvido, ES-C764-memorizar, ES-C764-revivir, ES-C764-diario

\chapter{Recuerdos}
\label{ch:es-recuerdos}
\hlchaptermodality{764}

\begin{chapteropening}
\textbf{By the end of this chapter:} \emph{I can talk about unforgettable moments, forgetting things and learning by heart.}

Complete the last lesson of chapter 764: I can talk about unforgettable moments, forgetting things and learning by heart.
\end{chapteropening}

\section[inolvidable]{inolvidable — unforgettable}

\label{lesson:ES-C764-inolvidable}



\begin{quote}
Before the new one: say the Spanish for the painter, then the Spanish for to photograph.
\end{quote}

\subsection*{You'll want to know: inolvidable}
\textbf{inolvidable} — \textquotedblleft{}unforgettable\textquotedblright{}.

\begin{quote}
Fue un viaje inolvidable.
\end{quote}

\begin{quote}
\emph{It was an unforgettable trip.}
\end{quote}

\begin{cousinweb}
\emph{inolvidable} is built on \emph{olvidar} \textquotedblleft{}to forget\textquotedblright{}, with \emph{in-} and \emph{-able}; \emph{olvidar} is inherited from a spoken Latin \emph{oblītāre} (EtymDB).
\end{cousinweb}

\begin{grammarlens}[title={What you've built}]
One more word for telling, explaining and describing at B1.
\end{grammarlens}

\subsection*{Your turn}
\begin{itemize}
\raggedright
  \item \emph{Say it:} \emph{inolvidable}
  \item \emph{Say it:} \emph{inolvidable}, once more
  \item \emph{Say it:} a sentence of your own with \emph{inolvidable}
  \item \emph{Recall:} say the Spanish for the painter, then the Spanish for to photograph, then say \emph{inolvidable} again
\end{itemize}

\begin{tcolorbox}[breakable,colback=teal!4,colframe=teal!35!black,title={Before you move on}]
What does inolvidable mean? (\textquotedblleft{}Unforgettable\textquotedblright{}.) Say it once more.
\end{tcolorbox}
\section[el olvido]{el olvido — forgetting; oblivion; oversight}

\label{lesson:ES-C764-olvido}



\begin{quote}
Before the new one: say the Spanish for to photograph, then the Spanish for unforgettable.
\end{quote}

\subsection*{You'll want to know: el olvido}
\textbf{el olvido} — \textquotedblleft{}forgetting; oblivion; oversight\textquotedblright{}.

\begin{quote}
Perdón, fue un olvido.
\end{quote}

\begin{quote}
\emph{Sorry, I forgot.}
\end{quote}

\emph{Caer en el olvido} — to be forgotten.

\begin{cousinweb}
\emph{olvido} is built on \emph{olvidar}, which is inherited from a spoken Latin \emph{oblītāre} (EtymDB), from \emph{oblītus}, the participle of \emph{oblīvīscī} \textquotedblleft{}to forget\textquotedblright{}.
\end{cousinweb}

\begin{grammarlens}[title={What you've built}]
One more word for telling, explaining and describing at B1.
\end{grammarlens}

\subsection*{Your turn}
\begin{itemize}
\raggedright
  \item \emph{Say it:} \emph{el olvido}
  \item \emph{Say it:} \emph{el olvido}, once more
  \item \emph{Say it:} a sentence of your own with \emph{el olvido}
  \item \emph{Recall:} say the Spanish for to photograph, then the Spanish for unforgettable, then say \emph{el olvido} again
\end{itemize}

\begin{tcolorbox}[breakable,colback=teal!4,colframe=teal!35!black,title={Before you move on}]
What does el olvido mean? (\textquotedblleft{}Forgetting; oblivion; oversight\textquotedblright{}.) Say it once more.
\end{tcolorbox}
\section[memorizar]{memorizar — to memorize, to learn by heart}

\label{lesson:ES-C764-memorizar}



\begin{quote}
Before the new one: say the Spanish for unforgettable, then the Spanish for forgetting.
\end{quote}

\subsection*{You'll want to know: memorizar}
\textbf{memorizar} — \textquotedblleft{}to memorize, to learn by heart\textquotedblright{}.

\begin{quote}
Tengo que memorizar estos números.
\end{quote}

\begin{quote}
\emph{I have to learn these numbers by heart.}
\end{quote}

\emph{Aprender de memoria} means the same.

\begin{cousinweb}
\emph{memorizar} is built on \emph{memoria} \textquotedblleft{}memory\textquotedblright{} with the ending \emph{-izar}; \emph{memoria} is a learned borrowing from Latin \emph{memoria}, from \emph{memor} \textquotedblleft{}mindful\textquotedblright{}.
\end{cousinweb}

\begin{grammarlens}[title={What you've built}]
One more word for telling, explaining and describing at B1.
\end{grammarlens}

\subsection*{Your turn}
\begin{itemize}
\raggedright
  \item \emph{Say it:} \emph{memorizar}
  \item \emph{Say it:} \emph{memorizar}, once more
  \item \emph{Say it:} a sentence of your own with \emph{memorizar}
  \item \emph{Recall:} say the Spanish for unforgettable, then the Spanish for forgetting, then say \emph{memorizar} again
\end{itemize}

\begin{tcolorbox}[breakable,colback=teal!4,colframe=teal!35!black,title={Before you move on}]
What does memorizar mean? (\textquotedblleft{}To memorize, to learn by heart\textquotedblright{}.) Say it once more.
\end{tcolorbox}
\section[revivir]{revivir — to relive; to revive}

\label{lesson:ES-C764-revivir}



\begin{quote}
Before the new one: say the Spanish for forgetting, then the Spanish for to memorize.
\end{quote}

\subsection*{You'll want to know: revivir}
\textbf{revivir} — \textquotedblleft{}to relive; to revive\textquotedblright{}.

\begin{quote}
Al ver las fotos revivo aquel día.
\end{quote}

\begin{quote}
\emph{Looking at the photos, I relive that day.}
\end{quote}

\begin{cousinweb}
\emph{revivir} is usually traced to Latin \emph{revīvere} \textquotedblleft{}to live again\textquotedblright{}, built on \emph{vīvere} \textquotedblleft{}to live\textquotedblright{}.
\end{cousinweb}

\begin{grammarlens}[title={What you've built}]
One more word for telling, explaining and describing at B1.
\end{grammarlens}

\subsection*{Your turn}
\begin{itemize}
\raggedright
  \item \emph{Say it:} \emph{revivir}
  \item \emph{Say it:} \emph{revivir}, once more
  \item \emph{Say it:} a sentence of your own with \emph{revivir}
  \item \emph{Recall:} say the Spanish for forgetting, then the Spanish for to memorize, then say \emph{revivir} again
\end{itemize}

\begin{tcolorbox}[breakable,colback=teal!4,colframe=teal!35!black,title={Before you move on}]
What does revivir mean? (\textquotedblleft{}To relive; to revive\textquotedblright{}.) Say it once more.
\end{tcolorbox}
\section[el diario]{el diario — the diary; the daily paper}

\label{lesson:ES-C764-diario}



\begin{quote}
Before the new one: say the Spanish for to memorize, then the Spanish for to relive.
\end{quote}

\subsection*{You'll want to know: el diario}
\textbf{el diario} — \textquotedblleft{}the diary; the daily paper\textquotedblright{}.

\begin{quote}
De pequeña escribía un diario.
\end{quote}

\begin{quote}
\emph{As a girl I kept a diary.}
\end{quote}

Also a daily newspaper, and the adjective daily: \emph{la vida diaria}. \emph{A diario} — every day.

\begin{cousinweb}
\emph{diario} comes from Latin \emph{diārium} \textquotedblleft{}a daily allowance\textquotedblright{} (EtymDB), built on \emph{diēs} \textquotedblleft{}day\textquotedblright{}.
\end{cousinweb}

\begin{grammarlens}[title={What you've built}]
One more word for telling, explaining and describing at B1.
\end{grammarlens}

\subsection*{Your turn}
\begin{itemize}
\raggedright
  \item \emph{Say it:} \emph{el diario}
  \item \emph{Say it:} \emph{el diario}, once more
  \item \emph{Say it:} a sentence of your own with \emph{el diario}
  \item \emph{Recall:} say the Spanish for to memorize, then the Spanish for to relive, then say \emph{el diario} again
\end{itemize}

\begin{tcolorbox}[breakable,colback=teal!4,colframe=teal!35!black,title={Before you move on}]
What does el diario mean? (\textquotedblleft{}The diary; the daily paper\textquotedblright{}.) Say it once more.
\end{tcolorbox}
